\documentclass[12pt]{article}
\usepackage[UTF8]{inputenc}
\usepackage{thaisupp}
\usepackage{enumitem}
\begin{document}
\title{ไฟฟ้าเหนี่ยวนำ}
\maketitle
\section*{In-class Questions}
\begin{enumerate}[label=\arabic*.]
\item กฎการเหนี่ยวนำของฟาราเดย์ (Faraday's law) กล่าวว่าอย่างไร \hfill \textbf{\small (Main)}
\begin{enumerate}[label=)]
\item ก) แรงเคลื่อนไฟฟ้าเหนี่ยวนำ (emf) แปรผกผันกับอัตราการเปลี่ยนแปลงของฟลักซ์แม่เหล็ก
\item ข) แรงเคลื่อนไฟฟ้าเหนี่ยวนำ (emf) แปรผันตรงกับอัตราการเปลี่ยนแปลงของฟลักซ์แม่เหล็ก
\item ค) แรงเคลื่อนไฟฟ้าเหนี่ยวนำมีค่าคงที่เสมอ
\item ง) แรงเคลื่อนไฟฟ้าเหนี่ยวนำขึ้นกับทิศทางของสนามแม่เหล็กเท่านั้น
\end{enumerate}
\item สูตรของกฎการเหนี่ยวนำของฟาราเดย์คือข้อใด \hfill \textbf{\small (Main)}
\begin{enumerate}[label=)]
\item ก) ℰ = dΦ/dt
\item ข) ℰ = -dΦ/dt
\item ค) ℰ = Φ·dI/dt
\item ง) ℰ = I·dΦ/dt
\end{enumerate}
\item กฎเลนซ์ (Lenz's law) กล่าวว่าอย่างไร \hfill \textbf{\small (Main)}
\begin{enumerate}[label=)]
\item ก) ทิศทางของกระแสไฟฟ้าเหนี่ยวนำจะต้านการเปลี่ยนแปลงของฟลักซ์แม่เหล็กที่ทำให้เกิดมันขึ้น
\item ข) ทิศทางของกระแสไฟฟ้าเหนี่ยวนำจะเสริมการเปลี่ยนแปลงของฟลักซ์แม่เหล็กที่ทำให้เกิดมันขึ้น
\item ค) ทิศทางของกระแสไฟฟ้าเหนี่ยวนำจะขนานกับสนามแม่เหล็กเสมอ
\item ง) ทิศทางของกระแสไฟฟ้าเหนี่ยวนำจะตั้งฉากกับสนามแม่เหล็กเสมอ
\end{enumerate}
\item วิธีใดต่อไปนี้ไม่สามารถทำให้เกิดกระแสไฟฟ้าเหนี่ยวนำในขดลวดได้ \hfill \textbf{\small (Main)}
\begin{enumerate}[label=)]
\item ก) เคลื่อนที่แม่เหล็กเข้าใกล้ขดลวด
\item ข) หมุนขดลวดในสนามแม่เหล็ก
\item ค) วางขดลวดนิ่งในสนามแม่เหล็กคงที่
\item ง) เปลี่ยนพื้นที่ของขดลวดในสนามแม่เหล็กคงที่
\end{enumerate}
\item เมื่อนำแม่เหล็กเข้าใกล้ขดลวดปิดอย่างรวดเร็ว กระแสไฟฟ้าเหนี่ยวนำจะมีทิศทางอย่างไร \hfill \textbf{\small (Main)}
\begin{enumerate}[label=)]
\item ก) ไหลในทิศทางที่ทำให้ขั้วใกล้แม่เหล็กเป็นขั้วเดียวกันกับแม่เหล็ก
\item ข) ไหลในทิศทางที่ทำให้ขั้วใกล้แม่เหล็กเป็นขั้วต่างชนิดกับแม่เหล็ก
\item ค) ไม่เกิดกระแสเหนี่ยวนำ
\item ง) ไหลกลับทิศทางตลอดเวลา
\end{enumerate}
\item ความเหนี่ยวนำตนเอง (self-inductance) ของขดลวดคืออะไร \hfill \textbf{\small (Main)}
\begin{enumerate}[label=)]
\item ก) ความสามารถของขดลวดในการเก็บประจุ
\item ข) ความสามารถของขดลวดในการต้านการเปลี่ยนแปลงของกระแส
\item ค) ความสามารถของขดลวดในการนำไฟฟ้า
\item ง) ความสามารถของขดลวดในการเปลี่ยนพลังงานไฟฟ้าเป็นพลังงานความร้อน
\end{enumerate}
\item หน่วยของความเหนี่ยวนำ (L) ในระบบ SI คืออะไร \hfill \textbf{\small (Main)}
\begin{enumerate}[label=)]
\item ก) โอห์ม (Ω)
\item ข) เฮนรี่ (H)
\item ค) เทสลา (T)
\item ง) เวเบอร์ (Wb)
\end{enumerate}
\item พลังงานที่เก็บสะสมในขดลวดที่มีความเหนี่ยวนำ L และมีกระแส I ไหลผ่าน มีค่าเท่าใด \hfill \textbf{\small (Main)}
\begin{enumerate}[label=)]
\item ก) U = 1/2 LI²
\item ข) U = LI²
\item ค) U = 1/2 LI
\item ง) U = L²I
\end{enumerate}
\item ขดลวด 100 รอบ มีฟลักซ์แม่เหล็กเปลี่ยนจาก 0.01 Wb เป็น 0.05 Wb ในเวลา 0.02 วินาที emf เหนี่ยวนำมีค่าเท่าใด \hfill \textbf{\small (Main)}
\begin{enumerate}[label=)]
\item ก) 100 V
\item ข) 200 V
\item ค) 50 V
\item ง) 250 V
\end{enumerate}
\item เมื่อนำแม่เหล็กออกจากขดลวดปิด กระแสเหนี่ยวนำจะไหลในทิศทางที่ทำให้เกิดอะไร \hfill \textbf{\small (Main)}
\begin{enumerate}[label=)]
\item ก) ทำให้ฟลักซ์เพิ่มขึ้นเพื่อต้านการลดลง
\item ข) ทำให้ฟลักซ์ลดลงเพื่อเสริมการลดลง
\item ค) ทำให้ฟลักซ์เพิ่มขึ้นเพื่อเสริมการลดลง
\item ง) ไม่มีผลต่อฟลักซ์
\end{enumerate}
\item ความเหนี่ยวนำตนเอง L ของขดลวดโซลินอยด์ขึ้นอยู่กับปริมาณใด \hfill \textbf{\small (Main)}
\begin{enumerate}[label=)]
\item ก) จำนวนรอบ พื้นที่หน้าตัด และความยาว
\item ข) กระแสไฟฟ้าที่ไหลผ่านเท่านั้น
\item ค) ความต้านทานของขดลวดเท่านั้น
\item ง) สนามแม่เหล็กภายนอกเท่านั้น
\end{enumerate}
\item ขดลวดมีความเหนี่ยวนำ 0.5 H มีกระแส 4 A ไหลผ่าน พลังงานที่เก็บสะสมในขดลวดมีค่าเท่าใด \hfill \textbf{\small (Main)}
\begin{enumerate}[label=)]
\item ก) 2 J
\item ข) 4 J
\item ค) 8 J
\item ง) 1 J
\end{enumerate}
\item ถ้าฟลักซ์แม่เหล็กที่ผ่านขดลวดเพิ่มขึ้นเป็น 3 เท่าในเวลาเท่าเดิม emf เหนี่ยวนำจะเปลี่ยนแปลงอย่างไร \hfill \textbf{\small (Main)}
\begin{enumerate}[label=)]
\item ก) เพิ่มเป็น 3 เท่า
\item ข) เพิ่มเป็น 9 เท่า
\item ค) เพิ่มเป็น √3 เท่า
\item ง) ไม่เปลี่ยนแปลง
\end{enumerate}
\item ขดลวด 200 รอบ มีพื้นที่ 0.01 ตารางเมตร หมุนในสนามแม่เหล็ก 0.5 T จากตำแหน่งที่ระนาบขดลวดตั้งฉากกับสนามเป็นตำแหน่งที่ระนาบขดลวดขนานกับสนาม การเปลี่ยนแปลงฟลักษ์มีค่าเท่าใด \hfill \textbf{\small (Main)}
\begin{enumerate}[label=)]
\item ก) 0.005 Wb
\item ข) 0.05 Wb
\item ค) 0.5 Wb
\item ง) 5 × 10⁻⁴ Wb
\end{enumerate}
\item ขดลวดโซลินอยด์ยาว 0.5 เมตร มีพื้นที่หน้าตัด 2 × 10⁻⁴ ตารางเมตร มีจำนวนรอบ 500 รอบ ความเหนี่ยวนำมีค่าเท่าใด (μ₀ = 4π × 10⁻⁷ H/m) \hfill \textbf{\small (Main)}
\begin{enumerate}[label=)]
\item ก) 1.26 × 10⁻⁴ H
\item ข) 1.26 × 10⁻⁵ H
\item ค) 1.26 × 10⁻³ H
\item ง) 1.26 × 10⁻⁶ H
\end{enumerate}
\item เมื่อเพิ่มความเร็วในการนำแม่เหล็กเข้าใกล้ขดลวดเป็น 2 เท่า emf เหนี่ยวนำจะเปลี่ยนแปลงอย่างไร \hfill \textbf{\small (Main)}
\begin{enumerate}[label=)]
\item ก) เพิ่มเป็น 2 เท่า
\item ข) เพิ่มเป็น 4 เท่า
\item ค) ลดลงเป็น 1/2 เท่า
\item ง) ไม่เปลี่ยนแปลง
\end{enumerate}
\item ขดลวดมีความเหนี่ยวนำ 0.2 H มีกระแสไหลเปลี่ยนแปลงจาก 5 A เป็น 3 A ในเวลา 0.1 วินาที จงหาพลังงานที่ปล่อยออกจากขดลวด \hfill \textbf{\small (Main)}
\begin{enumerate}[label=)]
\item ก) 1.6 J
\item ข) 3.2 J
\item ค) 0.8 J
\item ง) 0.4 J
\end{enumerate}
\item ขดลวด 50 รอบ มีพื้นที่ 0.02 ตารางเมตร วางในสนามแม่เหล็กที่เปลี่ยนแปลงจาก 0.1 T เป็น 0.6 T ใน 0.05 วินาที โดยสนามตั้งฉากกับระนาบขดลวด emf เหนี่ยวนำมีค่าเท่าใด \hfill \textbf{\small (Main)}
\begin{enumerate}[label=)]
\item ก) 10 V
\item ข) 100 V
\item ค) 1 V
\item ง) 50 V
\end{enumerate}
\item ถ้าเพิ่มจำนวนรอบของขดลวดโซลินอยด์เป็น 2 เท่า ความเหนี่ยวนำจะเปลี่ยนแปลงอย่างไร \hfill \textbf{\small (Main)}
\begin{enumerate}[label=)]
\item ก) เพิ่มเป็น 2 เท่า
\item ข) เพิ่มเป็น 4 เท่า
\item ค) ลดลงเป็น 1/2 เท่า
\item ง) เพิ่มเป็น √2 เท่า
\end{enumerate}
\item เมื่อหมุนขดลวดในสนามแม่เหล็กคงที่ ข้อใดเกิดขึ้น \hfill \textbf{\small (Main)}
\begin{enumerate}[label=)]
\item ก) เกิด emf เหนี่ยวนำคงที่
\item ข) เกิด emf เหนี่ยวนำที่เปลี่ยนแปลงเป็นฟังก์ชันไซน์ของเวลา
\item ค) ไม่เกิด emf เหนี่ยวนำ
\item ง) เกิดกระแสคงที่
\end{enumerate}
\item ความหนาแน่นพลังงานในสนามแม่เหล็ก (u = B²/2μ₀) มีหน่วยเป็นอะไร \hfill \textbf{\small (Main)}
\begin{enumerate}[label=)]
\item ก) J/m
\item ข) J/m²
\item ค) J/m³
\item ง) W/m³
\end{enumerate}
\item เมื่อปิดสวิตช์ในวงจรที่มีขดลวดเหนี่ยวนำ กระแสจะเพิ่มขึ้นอย่างรวดเร็วหรือช้า เพราะอะไร \hfill \textbf{\small (Main)}
\begin{enumerate}[label=)]
\item ก) เพิ่มขึ้นช้า เพราะ emf เหนี่ยวนำต้านการเพิ่มขึ้นของกระแส
\item ข) เพิ่มขึ้นเร็ว เพราะไม่มี emf เหนี่ยวนำ
\item ค) เพิ่มขึ้นช้า เพราะความต้านทานของขดลวดสูง
\item ง) เพิ่มขึ้นทันที เพราะขดลวดมีความเหนี่ยวนำต่ำ
\end{enumerate}
\item เมื่อนำแม่เหล็กวางนิ่งแล้วเลื่อนขดลวดออกมา แรงที่ต้องใช้เลื่อนขดลวดมีทิศทางอย่างไร \hfill \textbf{\small (Main)}
\begin{enumerate}[label=)]
\item ก) ตามทิศทางการเลื่อนขดลวด
\item ข) ตรงข้ามทิศทางการเลื่อนขดลวด
\item ค) ตั้งฉากกับทิศทางการเลื่อนขดลวด
\item ง) เป็นศูนย์
\end{enumerate}
\item ถ้าใส่แกนเหล็กอ่อนเข้าไปในขดลวดโซลินอยด์ ความเหนี่ยวนำจะเปลี่ยนแปลงอย่างไร \hfill \textbf{\small (Main)}
\begin{enumerate}[label=)]
\item ก) เพิ่มขึ้นมาก
\item ข) ลดลงมาก
\item ค) ไม่เปลี่ยนแปลง
\item ง) เพิ่มขึ้นเล็กน้อย
\end{enumerate}
\item ขดลวดวงแหวนรัศมี 0.1 เมตร มี 100 รอบ วางในสนามแม่เหล็กที่เปลี่ยนแปลงด้วยอัตรา 0.5 T/s จงหา emf เหนี่ยวนำในขดลวด \hfill \textbf{\small (Main)}
\begin{enumerate}[label=)]
\item ก) 1.57 V
\item ข) 0.157 V
\item ค) 15.7 V
\item ง) 0.0157 V
\end{enumerate}
\item การเปลี่ยนแปลงของปริมาณใดที่ทำให้เกิด emf เหนี่ยวนำ \hfill \textbf{\small (Extra)}
\begin{enumerate}[label=)]
\item ก) การเปลี่ยนแปลงของกระแสไฟฟ้า
\item ข) การเปลี่ยนแปลงของฟลักซ์แม่เหล็ก
\item ค) การเปลี่ยนแปลงของความต้านทาน
\item ง) การเปลี่ยนแปลงของแรงดันไฟฟ้า
\end{enumerate}
\item เมื่อนำแม่เหล็กเข้าใกล้ขดลวดปิดแล้วหยุดนิ่ง กระแสเหนี่ยวนำจะเป็นอย่างไร \hfill \textbf{\small (Extra)}
\begin{enumerate}[label=)]
\item ก) ไหลต่อเนื่อง
\item ข) ไหลชั่วขณะแล้วหยุด
\item ค) เป็นศูนย์
\item ง) ไหลกลับทิศทาง
\end{enumerate}
\item เครื่องหมายลบในสูตร ℰ = -N dΦ/dt มาจากอะไร \hfill \textbf{\small (Extra)}
\begin{enumerate}[label=)]
\item ก) กฎของนิวตัน
\item ข) กฎเลนซ์
\item ค) กฎของฟาราเดย์
\item ง) กฎของคูลอมบ์
\end{enumerate}
\item ขดลวดที่มีความเหนี่ยวนำมาก จะต้านการเปลี่ยนแปลงของกระแสได้มากหรือน้อย \hfill \textbf{\small (Extra)}
\begin{enumerate}[label=)]
\item ก) มาก
\item ข) น้อย
\item ค) เท่ากัน
\item ง) ขึ้นกับความต้านทาน
\end{enumerate}
\item พลังงานที่เก็บสะสมในขดลวดเหนี่ยวนำอยู่ในรูปอะไร \hfill \textbf{\small (Extra)}
\begin{enumerate}[label=)]
\item ก) พลังงานความร้อน
\item ข) พลังงานสนามแม่เหล็ก
\item ค) พลังงานสนามไฟฟ้า
\item ง) พลังงานศักย์ไฟฟ้า
\end{enumerate}
\item เมื่อนำแม่เหล็กเข้าใกล้ขดลวดปิด แรงที่เกิดขึ้นระหว่างแม่เหล็กกับขดลวดจะเป็นอย่างไร \hfill \textbf{\small (Extra)}
\begin{enumerate}[label=)]
\item ก) เป็นแรงดึงดูด
\item ข) เป็นแรงผลัก
\item ค) เป็นศูนย์
\item ง) เปลี่ยนทิศทางตลอดเวลา
\end{enumerate}
\item ขดลวด 500 รอบ มีฟลักซ์แม่เหล็กลดลงจาก 0.08 Wb เป็น 0.02 Wb ใน 0.03 วินาที emf เหนี่ยวนำมีค่าเท่าใด \hfill \textbf{\small (Extra)}
\begin{enumerate}[label=)]
\item ก) 1000 V
\item ข) 500 V
\item ค) 1500 V
\item ง) 2000 V
\end{enumerate}
\item ถ้าเพิ่มความยาวของขดลวดโซลินอยด์เป็น 2 เท่า โดยไม่เปลี่ยนจำนวนรอบต่อหน่วยความยาว ความเหนี่ยวนำจะเปลี่ยนแปลงอย่างไร \hfill \textbf{\small (Extra)}
\begin{enumerate}[label=)]
\item ก) เพิ่มเป็น 2 เท่า
\item ข) ลดลงเป็น 1/2 เท่า
\item ค) เพิ่มเป็น 4 เท่า
\item ง) ไม่เปลี่ยนแปลง
\end{enumerate}
\item กระแสเหนี่ยวนำในขดลวดปิดขึ้นกับอะไร \hfill \textbf{\small (Extra)}
\begin{enumerate}[label=)]
\item ก) emf เหนี่ยวนำและความต้านทานของขดลวด
\item ข) emf เหนี่ยวนำเท่านั้น
\item ค) ความต้านทานของขดลวดเท่านั้น
\item ง) จำนวนรอบของขดลวดเท่านั้น
\end{enumerate}
\item ขดลวดเหนี่ยวนำมีกระแส 10 A ไหลผ่าน พลังงานที่เก็บสะสมเป็น 50 J ความเหนี่ยวนำของขดลวดมีค่าเท่าใด \hfill \textbf{\small (Extra)}
\begin{enumerate}[label=)]
\item ก) 0.5 H
\item ข) 1 H
\item ค) 2 H
\item ง) 0.25 H
\end{enumerate}
\item เมื่อหมุนขดลวดในสนามแม่เหล็กด้วยความเร็วเชิงมุม ω emf สูงสุดจะเท่ากับข้อใด \hfill \textbf{\small (Extra)}
\begin{enumerate}[label=)]
\item ก) ℰ₀ = NBAω
\item ข) ℰ₀ = NBω
\item ค) ℰ₀ = NAω
\item ง) ℰ₀ = BAω
\end{enumerate}
\item เมื่อนำขดลวดวงแหวนเข้าใกล้แม่เหล็ก แรงผลักที่เกิดขึ้นระหว่างขดลวดกับแม่เหล็กเกิดจากอะไร \hfill \textbf{\small (Extra)}
\begin{enumerate}[label=)]
\item ก) แรงโน้มถ่วง
\item ข) กระแสเหนี่ยวนำที่เกิดในขดลวด
\item ค) สนามไฟฟ้าสถิต
\item ง) แรงระหว่างประจุบวกและประจุลบ
\end{enumerate}
\item ขดลวดโซลินอยด์ยาว 0.2 เมตร มีเส้นผ่านศูนย์กลาง 0.05 เมตร มีจำนวนรอบ 200 รอบ มีความเหนี่ยวนำเท่าใด (μ₀ = 4π × 10⁻⁷ H/m) \hfill \textbf{\small (Extra)}
\begin{enumerate}[label=)]
\item ก) 9.87 × 10⁻⁶ H
\item ข) 9.87 × 10⁻⁵ H
\item ค) 9.87 × 10⁻⁷ H
\item ง) 9.87 × 10⁻⁴ H
\end{enumerate}
\item ขดลวดมีความเหนี่ยวนำ 0.1 H มีกระแสเพิ่มจาก 0 A เป็น 5 A อย่างสม่ำเสมอใน 0.2 วินาที จงหา emf เหนี่ยวนำตนเอง \hfill \textbf{\small (Extra)}
\begin{enumerate}[label=)]
\item ก) 2.5 V
\item ข) 5 V
\item ค) 1.25 V
\item ง) 0.5 V
\end{enumerate}
\item เมื่อเปิดสวิตช์ในวงจรที่มีขดลวดเหนี่ยวนำ เกิดอะไรขึ้น \hfill \textbf{\small (Extra)}
\begin{enumerate}[label=)]
\item ก) กระแสลดลงทันที
\item ข) เกิดประกายไฟที่สวิตช์
\item ค) ไม่มีผลอะไร
\item ง) กระแสเพิ่มขึ้นชั่วขณะ
\end{enumerate}
\item ขดลวด 100 รอบ มีพื้นที่ 0.05 ตารางเมตร วางในสนามแม่เหล็ก 0.8 T ที่ทำมุม 60° กับระนาบขดลวด ฟลักษ์แม่เหล็กที่ผ่านขดลวดมีค่าเท่าใด \hfill \textbf{\small (Extra)}
\begin{enumerate}[label=)]
\item ก) 0.04 Wb
\item ข) 0.02 Wb
\item ค) 0.04√3 Wb
\item ง) 0.01 Wb
\end{enumerate}
\item ถ้าใส่แกนเหล็กอ่อนที่มีความ проницаемость μ = 1000μ₀ เข้าไปในขดลวดโซลินอยด์ ความเหนี่ยวนำจะเปลี่ยนแปลงอย่างไร \hfill \textbf{\small (Extra)}
\begin{enumerate}[label=)]
\item ก) เพิ่มเป็น 1000 เท่า
\item ข) เพิ่มเป็น 100 เท่า
\item ค) เพิ่มเป็น √1000 เท่า
\item ง) ไม่เปลี่ยนแปลง
\end{enumerate}
\item ขดลวดโซลินอยด์มีความหนาแน่นพลังงานในสนามแม่เหล็ก 100 J/m³ สนามแม่เหล็กภายในขดลวดมีค่าเท่าใด (μ₀ = 4π × 10⁻⁷ H/m) \hfill \textbf{\small (Extra)}
\begin{enumerate}[label=)]
\item ก) 0.5 T
\item ข) 0.5 × 10⁻³ T
\item ค) 0.5 × 10³ T
\item ง) 0.25 T
\end{enumerate}
\item ขดลวดวงแหวนรัศมี 0.1 เมตร มีความต้านทาน 0.02 Ω วางในสนามแม่เหล็กที่เปลี่ยนแปลง 2 T/s emf เหนี่ยวนำมีค่าเท่าใด และกระแสเหนี่ยวนำมีค่าเท่าใด \hfill \textbf{\small (Extra)}
\begin{enumerate}[label=)]
\item ก) 0.0628 V, 3.14 A
\item ข) 0.628 V, 31.4 A
\item ค) 6.28 V, 314 A
\item ง) 0.00628 V, 0.314 A
\end{enumerate}
\item เมื่อนำแม่เหล็กเข้าใกล้ขดลวดทองแดงที่หมุนอยู่ แรงที่ต้องใช้หมุนขดลวดเป็นอย่างไร \hfill \textbf{\small (Extra)}
\begin{enumerate}[label=)]
\item ก) น้อยกว่าปกติ
\item ข) มากกว่าปกติ
\item ค) เท่ากับปกติ
\item ง) เป็นศูนย์
\end{enumerate}
\item ถ้าต้องการเพิ่มความเหนี่ยวนำของขดลวดโซลินอยด์ 4 เท่า โดยไม่เปลี่ยนจำนวนรอบ ต้องเปลี่ยนแปลงอะไร \hfill \textbf{\small (Extra)}
\begin{enumerate}[label=)]
\item ก) เพิ่มความยาวเป็น 2 เท่า
\item ข) เพิ่มพื้นที่หน้าตัดเป็น 4 เท่า
\item ค) ลดพื้นที่หน้าตัดเป็น 1/4
\item ง) เพิ่มความยาวเป็น 4 เท่า
\end{enumerate}
\item ขดลวดเหนี่ยวนำสองขดลวดมีความเหนี่ยวนำ L₁ = 2 H และ L₂ = 8 H ถ้าต่อขดลวดทั้งสองอนุกรมกัน ความเหนี่ยวนำรวมเป็นเท่าใด \hfill \textbf{\small (Extra)}
\begin{enumerate}[label=)]
\item ก) 10 H
\item ข) 6 H
\item ค) 4 H
\item ง) 16 H
\end{enumerate}
\item เมื่อลวดตัวนำเคลื่อนที่ตัดเส้นแรงแม่เหล็กในสนามแม่เหล็ก B ด้วยความเร็ว v ในแนวตั้งฉากกับทั้งความเร็วและสนาม emf ที่เหนี่ยวนำมีค่าเท่าใด \hfill \textbf{\small (Extra)}
\begin{enumerate}[label=)]
\item ก) ℰ = Bvl
\item ข) ℰ = Bv/l
\item ค) ℰ = Bv²l
\item ง) ℰ = Bl/v
\end{enumerate}
\item ขดลวดเหนี่ยวนำสองขดลวดมีความเหนี่ยวนำ L₁ = 3 H และ L₂ = 6 H ถ้าต่อขดลวดทั้งสองขนานกัน ความเหนี่ยวนำรวมเป็นเท่าใด \hfill \textbf{\small (Extra)}
\begin{enumerate}[label=)]
\item ก) 2 H
\item ข) 9 H
\item ค) 3 H
\item ง) 1/2 H
\end{enumerate}
\item แท่งตัวนำยาว 0.5 เมตร เคลื่อนที่ด้วยความเร็ว 10 m/s ตั้งฉากกับสนามแม่เหล็ก 0.2 T emf ที่เหนี่ยวนำมีค่าเท่าใด \hfill \textbf{\small (Extra)}
\begin{enumerate}[label=)]
\item ก) 1 V
\item ข) 10 V
\item ค) 0.1 V
\item ง) 100 V
\end{enumerate}
\item ขดลวดโซลินอยด์ 1000 รอบ มีความเหนี่ยวนำ 0.005 H ถ้าต้องการให้ความเหนี่ยวนำเป็น 0.02 H ต้องเพิ่มจำนวนรอบเป็นเท่าใด \hfill \textbf{\small (Extra)}
\begin{enumerate}[label=)]
\item ก) 2000 รอบ
\item ข) 4000 รอบ
\item ค) 1000 รอบ
\item ง) 500 รอบ
\end{enumerate}
\item ขดลวดโซลินอยด์มีความเหนี่ยวนำ 0.01 H มีกระแส 2 A ไหลผ่าน จงหาความหนาแน่นพลังงานในสนามแม่เหล็กภายในขดลวด (ปริมาตรของโซลินอยด์ = 10⁻⁴ m³) \hfill \textbf{\small (Extra)}
\begin{enumerate}[label=)]
\item ก) 2 × 10⁴ J/m³
\item ข) 2 × 10² J/m³
\item ค) 2 × 10³ J/m³
\item ง) 2 × 10⁵ J/m³
\end{enumerate}
\item เมื่อนำแม่เหล็กที่หมุนอยู่มาใกล้ขดลวดปิด แรงที่ต้องใช้ทำให้แม่เหล็กหมุนจะเป็นอย่างไร \hfill \textbf{\small (Extra)}
\begin{enumerate}[label=)]
\item ก) น้อยกว่าปกติ
\item ข) มากขึ้นเพราะกระแสเหนี่ยวนำต้านการเคลื่อนที่
\item ค) เท่าเดิม
\item ง) เป็นศูนย์
\end{enumerate}
\item ขดลวดตัวนำยาวมากมีความเหนี่ยวนำตนเอง L ถ้าพับขดลวดครึ่งหนึ่งโดยให้กระแสไหลในทิศทางเดียวกัน ความเหนี่ยวนำจะเปลี่ยนแปลงอย่างไร \hfill \textbf{\small (Extra)}
\begin{enumerate}[label=)]
\item ก) เพิ่มขึ้น 4 เท่า
\item ข) เพิ่มขึ้น 2 เท่า
\item ค) ลดลง 1/4 เท่า
\item ง) เพิ่มขึ้นเล็กน้อย
\end{enumerate}
\item ขดลวดรูปสี่เหลี่ยมผืนผ้ากว้าง 0.1 m ยาว 0.2 m มี 50 รอบ หมุนในสนามแม่เหล็ก 0.5 T ด้วยความเร็วเชิงมุม 100 rad/s emf สูงสุดที่เกิดขึ้นมีค่าเท่าใด \hfill \textbf{\small (Extra)}
\begin{enumerate}[label=)]
\item ก) 50 V
\item ข) 100 V
\item ค) 5 V
\item ง) 500 V
\end{enumerate}
\end{enumerate}
\end{document}
